\cotefile{136}
\scanpage{48}
% @01–02, upper right, boxed heading
\fbox{\shortstack{Cohomologie des ens. $1/2$ simpliciaux\\et \unsure{résolutions comme $1/2$ simpliciales|résolution comme $1/2$ simpliciale|résolutions canon. $1/2$ simpliciales} ?}}

% @03, opening
Soit $k$ un \deleted{groupe abélien}\inserted{anneau}, $M$ un $k$-module,
% @04, functor
\[
 C^*\colon\widehat{\Delta}^{\,\circ}\longrightarrow
 \text{complexes de cochaines de $k$-modules}
\]
% @05, hypotheses
un foncteur satisfaisant les \unsure{suppos.}

% @06, condition a
a) $\forall\ X_*\in\operatorname{Ob}\widehat{\Delta},\quad H^*(X_*,M)\simeq H^*(C^*(X_*))$ (isom fonctoriel en $M$ -- au moins pour $H^0$)

% @07, condition b
b) $\forall n\geq 0,\quad X_*\longmapsto C^n(X_*)$ transforme les \unsure{$\varinjlim$ en $\varprojlim$}.

% @08, representability
Alors b) signifie que les $C^n$ sont représentables
% @09, representation
\[
 C^n(X_*)=\operatorname{Hom}(X_*,C_*^n).
\]
% @10–14, resolution and augmentation
a) implique, en faisant $X_*=\Delta^m$, que les $C_{\unsure{m|k|\bullet}}^*$ \unsure{sont des} \uline{résolutions} de $\uline{M}$, alors que l'augmentation $M\to C_m^*$ \unsure{soit} fonctorielle en $m$, i.e. $C_*^*$ est une résolution de $M_*=\unsure{S_*|\underline{k}}\otimes_k M\simeq M_{\unsure{s.s.}}$ \unsure{ou $M$}, ce qui explique bien que l'on a, \unsure{aux conditions de} fonctorialité \unsure{près}, $H^*(C^*(X_*))=H^*(X_*,M)$,
\unsure{des suites spectrales}. \illegible{2} pour $X_*$ \unsure{de la forme} $\Delta^m$.
% @14a, left margin
\margin{\unsure{Cela l'intègre au programme ?}}

% @15–18, converse across the fold
\deleted{En prenant} $X_*$ \unsure{``arbitraire''} on trouve que \unsure{$H^k(X_*,M)$} \illegible{3} $\unsure{H^0(C^*(X_*))}=M$ \unsure{pour $H$}. \unsure{Il est a priori} \unsure{faussement|fraîchement} fort que $C_*^\bullet$ \unsure{soit} une résolution de $M_*$) i.e. \unsure{tout se trouve} $X_*\to Z_*^i$.
% @19, two lifting diagrams
\diagram{\begin{tikzcd}[ampersand replacement=\&,column sep=small,row sep=small]
 C_*^{i-1}\arrow[r]\&Z_*^i\&[-1em]\text{(\,pour}\&C_*^0\arrow[r]\&Z_*^1\&\text{)}\\
 \& X_*\arrow[u]\&\&\&X_*\arrow[u]\&
\end{tikzcd}}
% @20–23, lifting and flasqueness
se relève en $C_*^{i-1}$ \unsure{--}. On a envie de prouver que a) pour des quelconques signifie que les $C_*^i$ sont flasques, (ce qui serait \unsure{évidemment} suffisant pour qu'on ait a)).
% @23–25, cancelled lines
\deleted{\unsure{ce qu'il y a là} \illegible{3} \unsure{ce qui} \unsure{est démontré} que \illegible{2} fournissant \unsure{alors} a) \unsure{en résulte}. \inserted{\unsure{(cf. question dessus)}} \unsure{appliquer a)} \illegible{1}. On aura}
% @26, cohomology row
\[
 H^1(X_*,C_*^0)\to H^0(X_*,Z_*^1)\to H^1(X_*,M)\to H^1(X_*,C_*^0)\xrightarrow{\unsure{\text{inj}}} H^1(X_*,Z_*^1).
\]
% @27–29, paragraph
Mais c'est \unsure{faux} déjà pour $M=0$. Il suffit, pour qu'on ait \unsure{b)|a)}, que la suite exacte \unsure{longue} des $C_*^{\unsure{0|n}}\to Z_*^{\unsure{1|i}}\to\cdots$ soit ``splittée'' \unsure{trivialement|bêtement} : \unsure{homotopie} !

% @30–31, observation
\uline{NB} Si on a $C_*^\bullet$ une résolution flasque de $M_*$, les $Z_*^n(C_*^\bullet)$ sont des rétractes des $K_*(M,n)$ \ldots
% archive sheet number: 47

